% =============================================================================
%  lec14.tex —— 第 14 讲 谱（底本 PDF 第 110–116 页）
% =============================================================================
\lecture{14}{谱}
\begin{center}\small\color{black!55}（底本 PDF 第 110–116 页）\end{center}

本讲简要介绍稳定同伦论与谱。

从本讲起，我们使用记号
\[
\Spc\defeq\Grpd_{\infty}.
\]

\section{稳定同伦群}

\nitem\label{L14-1.1} 设 $\Top_{\ast}$ 是带点空间的普通范畴。存在一个伴随
\[
\Sigma:\Top_{\ast}\rightleftarrows\Top_{\ast}:\Omega,
\]
其中
\begin{itemize}[leftmargin=2.2em,itemsep=0.1em]
\item 左伴随 $\Sigma$ 是（带点）\textbf{悬挂函子}，定义为
\[
\Sigma X\defeq S^{1}\wedge X\defeq (S^{1}\times X)/((\{\ast\}\times X)\cup(S^{1}\times\{\ast\})).
\]
\item 右伴随 $\Omega$ 是\textbf{环路函子}，定义为
\[
\Omega Y\defeq\Hom_{\Top_{\ast}}(S^{1},Y),
\]
其中右端带有紧开拓扑。
\end{itemize}

\nitem\label{L14-1.2} 事实上，这个伴随与 Quillen 的经典模型结构相容\footnote{为此，我们必须把 $\Top$ 换成紧生成拓扑空间的范畴（以确保它是笛卡尔闭的）。任何 CW 复形都是紧生成的。}。取导出函子，得到伴随
\begin{equation}\label{L14-eq1.1}
L\Sigma:\hTop_{\ast}\rightleftarrows\hTop_{\ast}:R\Omega.
\end{equation}
由于带点 CW 复形都是双纤维的，我们有
\[
[\Sigma X,Y]\simeq[X,\Omega Y],
\]
其中 $[-,-]$ 表示连续映射的同伦类的集合。

\begin{exercise}\label{L14-1.3}
证明 $\Sigma S^{n}\simeq S^{n+1}$。
\end{exercise}

\begin{exercise}\label{L14-1.4}
对 $Y\in\Top_{\ast}$，存在典范同构 $\pi_{n+1}(Y)\simeq\pi_{n}(\Omega Y)$，其中右端的群结构由拼接映射 $\Omega Y\times\Omega Y\to\Omega Y$ 诱导。
\end{exercise}

\nitem\label{L14-1.5} 对带点 CW 复形 $X$ 与 $Y$，定义
\[
[X,Y]_{s}\defeq\colim_{k}[\Sigma^{k}X,\Sigma^{k}Y].
\]

\begin{exercise}\label{L14-1.6}
证明 $[X,Y]_{s}$ 自然地是一个阿贝尔群。提示：
\[
[\Sigma^{k+2}X,\Sigma^{k+2}Y]\simeq[\Sigma^{k}X,\Omega^{2}\Sigma^{k+2}Y].
\]
\end{exercise}

\begin{definition}\label{L14-1.7}
设 $Y$ 是带点 CW 复形。$Y$ 的\textbf{第 $n$ 个稳定同伦群}定义为
\[
\pi^{s}_{n}(Y)\defeq\colim_{k}\pi_{n+k}(\Sigma^{k}Y).
\]
\end{definition}

\begin{example}\label{L14-1.8}
群 $\pi_{n}(\mathbb{S})\defeq\pi^{s}_{n}(S^{0})$ 称为\textbf{球谱}的第 $n$ 个稳定同伦群。时至今日，人们已算到 $n\le 90$。
\end{example}

\nitem\label{L14-1.9} 稳定同伦论研究空间的稳定同伦群，更一般地，研究各种同伦不变量在悬挂函子 $\Sigma^{k}$ 下（$k\to\infty$ 时）的极限行为。相反，通常的同伦论被称为\textbf{非稳定同伦论}。我们的指导哲学是

\begin{slogan}\label{L14-1.10}\centering\par
稳定同伦论是非稳定同伦论的线性化：\\[0.4em]
稳定同伦论 $=$ 同伦论中的线性代数。
\end{slogan}

\section{谱}

\nitem\label{L14-2.1} 在前几讲中，我们已解释了下述哲学。为了捕捉 $\Top$ 中全部同伦不变的信息，我们需要的对象是空间构成的 $\infty$-范畴 $\Spc$，而不是它的同伦 $1$-范畴 $\mathrm{h}\Spc\simeq\hTop$。类似地，带点空间的同伦不变信息应由余切片 $\infty$-范畴
\[
\Spc_{\ast}\defeq\Spc_{\ast/}
\]
来捕捉。由此可见，稳定同伦论“正确”的舞台应当是 $\Spc_{\ast}$ 的 $\infty$-范畴\textbf{稳定化}或\textbf{线性化}。例如，我们希望对于（足够好的）对象 $X,Y\in\Spc_{\ast}$，这个稳定化 $\infty$-范畴中相应的映射空间由
\[
\colim_{k}\operatorname{Maps}_{\Spc_{\ast}}(\Sigma^{k}X,\Sigma^{k}Y)
\]
给出。让我们先定义 $\Sigma$ 与 $\Omega$ 的 $\infty$-范畴版本。

\begin{definition}\label{L14-2.2}
称 $\infty$-范畴 $\cC$ 是\textbf{带点的}，若它有一个对象 $0\in\cC$ 既是始对象又是终对象。我们称它为 $\cC$ 的\textbf{零对象}。
\end{definition}

\begin{exercise}\label{L14-2.3}
设 $\cC$ 是配有终对象 $\ast$ 的 $\infty$-范畴。证明 $\cC_{\ast/}$ 是带点的。特别地，$\Spc_{\ast}$ 是带点的。
\end{exercise}

\begin{definition}\label{L14-2.4}
设 $\cC$ 是配有有限余极限的带点 $\infty$-范畴。$\cC$ 上的\textbf{悬挂函子}定义为
\[
\Sigma:\cC\to\cC,\qquad X\mapsto 0\sqcup_{X}0.
\]
\end{definition}

\begin{definition}\label{L14-2.5}
设 $\cC$ 是配有有限极限的带点 $\infty$-范畴。$\cC$ 上的\textbf{环路函子}定义为
\[
\Omega:\cC\to\cC,\qquad Y\mapsto 0\times_{Y}0.
\]
\end{definition}

\begin{exercise}\label{L14-2.6}
设 $\cC$ 是同时配有有限极限与有限余极限的带点 $\infty$-范畴。构造一个伴随：
\[
\Sigma:\cC\rightleftarrows\cC:\Omega.
\]
\end{exercise}

\begin{exercise}\label{L14-2.7}
对 $\cC\defeq\Spc_{\ast}$，上述伴随诱导同伦范畴之间的伴随
\[
\mathrm{h}\Sigma:\mathrm{h}\Spc_{\ast}\rightleftarrows\mathrm{h}\Spc_{\ast}:\mathrm{h}\Omega.
\]
证明这个伴随经由等价 $\mathrm{h}\Spc_{\ast}\simeq\hTop_{\ast}$ 可等同于 \eqref{L14-eq1.1}。
\end{exercise}

\nitem\label{L14-2.8} 构造
\[
\operatorname{Maps}(-,-)\mapsto\colim_{k}\operatorname{Maps}(\Sigma^{k}(-),\Sigma^{k}(-))
\]
可以看作对函子 $\Sigma$ 作形式反演。

\begin{exercise}\label{L14-2.9}
设 $A$ 是交换环，$f\in A$ 是元素。证明
\[
A_{f}\simeq\colim[A\xrightarrow{f}A\xrightarrow{f}\cdots].
\]
\end{exercise}

\nitem\label{L14-2.10} 设 $\cC$ 是同时配有有限极限与有限余极限的带点 $\infty$-范畴。受上述构造启发，我们希望把 $\cC$ 的\textbf{稳定化}定义为
\[
\colim[\cC\xrightarrow{\Sigma}\cC\xrightarrow{\Sigma}\cdots].
\]
然而，我们必须小心这个余极限是在哪里取的。例如，当 $\cC$ 是可表现范畴（如 $\Spc_{\ast}$）时，我们希望得到一个可表现 $\infty$-范畴。

\begin{exercise}\label{L14-2.11}
设 $\cC$ 是带点的可表现 $\infty$-范畴。证明余极限
\[
\colim[\cC\xrightarrow{\Sigma}\cC\xrightarrow{\Sigma}\cdots]\in\mathrm{Pr}^{L}
\]
经由 $\mathrm{Pr}^{L}\simeq(\mathrm{Pr}^{R})^{\op}$ 对应对极限
\[
\lim[\cC\xleftarrow{\Omega}\cC\xleftarrow{\Omega}\cdots]\in\mathrm{Pr}^{R}.
\]
\end{exercise}

\nitem\label{L14-2.12} 回忆 $\mathrm{Pr}^{R}$ 中的极限可以按 $\widehat{\Cat}_{\infty}$ 中的极限来计算。这引向下述定义。

\begin{definition}\label{L14-2.13}
设 $\cC$ 是配有有限极限的带点 $\infty$-范畴。定义
\[
\mathrm{Sptr}(\cC)\defeq\lim[\cC\xleftarrow{\Omega}\cC\xleftarrow{\Omega}\cdots],
\]
并称之为 $\cC$ 的\textbf{谱对象} $\infty$-范畴。我们把对第 $(k+1)$ 项的求值态射记为
\[
\Omega^{\infty-k}:\mathrm{Sptr}(\cC)\to\cC.
\]
\end{definition}

\yizhu{本讲把 $\mathrm{Sptr}(\cC)$ 译作“$\cC$ 的谱对象 $\infty$-范畴”，把 $\mathrm{Sptr}(\Spc_{\ast})$ 简称为“谱 $\infty$-范畴”（例~\ref{L14-2.14}）；这里的“谱”与第 15 讲的“稳定 $\infty$-范畴”是不同层次的概念。\textbf{[译者按（无外部依据）]}}

\begin{example}\label{L14-2.14}
对 $\cC\defeq\Spc_{\ast}$，记
\[
\mathrm{Sptr}\defeq\mathrm{Sptr}(\Spc_{\ast}),
\]
并称之为\textbf{谱} $\infty$-范畴。
\end{example}

\begin{exercise}\label{L14-2.15}
证明 $\Omega:\cC\to\cC$ 保持有限极限。由此推出 $\mathrm{Sptr}(\cC)$ 配有有限极限，且函子 $\Omega^{\infty-k}$ 保持并检测有限极限。
\end{exercise}

\begin{exercise}\label{L14-2.16}
证明 $\mathrm{Sptr}(\cC)$ 是带点的。
\end{exercise}

\begin{exercise}\label{L14-2.17}
设 $\Omega_{\mathrm{Sptr}(\cC)}$ 是 $\mathrm{Sptr}(\cC)$ 上的环路函子。证明
\[
\Omega^{\infty-k}\circ\Omega_{\mathrm{Sptr}(\cC)}(E)\simeq\Omega^{\infty-k+1}(E).
\]
由此推出 $\Omega_{\mathrm{Sptr}(\cC)}$ 是等价。提示：
\begin{center}
\begin{tikzcd}[column sep=3.4em, row sep=3em]
\cC \arrow[d, "\Omega"] & \cC \arrow[l, "\Omega"'] \arrow[d, "\Omega"] & \cdots \arrow[l] \\
\cC & \cC \arrow[l, "\Omega"'] & \cdots \arrow[l]
\end{tikzcd}
\end{center}
\end{exercise}

\begin{exercise}\label{L14-2.18}
证明
\[
\Omega^{\infty-k}:\mathrm{Sptr}(\mathrm{Sptr}(\cC))\to\mathrm{Sptr}(\cC)
\]
是等价。
\end{exercise}

\begin{remark}\label{L14-2.19}
在下一讲中，我们将定义并研究\textbf{稳定 $\infty$-范畴}，它们恰是那些配有有限极限且 $\Omega$ 为等价的带点 $\infty$-范畴。
\end{remark}

\begin{exercise}\label{L14-2.20}
证明 $\mathrm{h}\mathrm{Sptr}(\cC)$ 是加法范畴。提示：
\[
\operatorname{Maps}_{\mathrm{Sptr}(\cC)}(E,E')\simeq\Omega^{2}\operatorname{Maps}_{\mathrm{Sptr}(\cC)}(E,\Sigma^{2}E').
\]
\end{exercise}

\section{谱与无穷环路空间}

\nitem\label{L14-3.1} 不严格地说，知道一个对象 $X\in\mathrm{Sptr}(\cC)$，等价于知道下述数据：
\begin{itemize}[leftmargin=2.2em,itemsep=0.1em]
\item 对任意 $n\ge 0$，一个对象 $X_{n}\in\cC$；
\item 对任意 $n\ge 0$，一个等价 $X_{n}\simeq\Omega X_{n+1}$。
\end{itemize}
这里我们取 $X_{n}$ 为 $\Omega^{\infty-k}X$。

\yizhu{原文此处作 $\Omega^{\infty-k}X$；按上下文，$X_{n}$ 是第 $n$ 项，故应为 $\Omega^{\infty-n}X$。译文照录原文，以便与底本对勘。\textbf{[译者按（无外部依据）]}}

注意，$X_{n+1}$ 连同等价 $X_{n}\simeq\Omega X_{n+1}$ 给出了 $X_{n}$ 的一个\textbf{去环路}（delooping）。于是我们得到下述口号。

\begin{slogan}\label{L14-3.2}\centering\par
谱就是配备无穷多次去环路结构的空间。
\end{slogan}

\begin{warning}\label{L14-3.3}
对空间 $Y\in\Spc_{\ast}$，它的去环路即便在同伦意义下也不唯一。因此在上面的数据中，记住所有的去环路至关重要。
\end{warning}

\nitem\label{L14-3.4} 注意，环路空间 $\Omega Z$ 带有同伦连贯的乘法结构，它使 $\pi_{0}(\Omega Z)$ 成为一个抽象群。在后续各讲中，我们将严格定义这种结构，并称之为\textbf{群样 $E_{1}$-结构}（grouplike $E_1$-structure）。此外，给定群样 $E_{1}$-空间 $Y$，存在本质唯一的连通去环路 $BY$，使得 $Y\simeq\Omega BY$ 且该等价与群样 $E_{1}$-结构相容。

此外，我们还将把上述讨论推广到迭代环路空间 $\Omega^{n}Z$ 与群样 $E_{n}$-空间。事实上，这对 $n=\infty$ 也成立，而我们将解释下述口号。

\begin{slogan}\label{L14-3.5}\centering\par
连接谱\footnote{称谱 $E\in\mathrm{Sptr}$ 是\textbf{连接的}，若 $n<0$ 时 $\pi_{n}E\simeq 0$。见下文定义~\ref{L14-4.7}。}就是群样 $E_{\infty}$-空间。
\end{slogan}

\yizhu{连接谱与群样 $E_{\infty}$-空间的等价是经典无穷环路空间理论的核心结论；参见 May, \emph{The Geometry of Iterated Loop Spaces}（1972），以及 Segal 关于无穷环路空间的工作。\textbf{[来源：A]}}

\section{空间与谱}

\nitem\label{L14-4.1} 本节聚焦 $\cC$ 带点且可表现的情形，例如 $\cC\defeq\Spc_{\ast}$。按定义，我们有余极限图
\[
[\cC\xrightarrow{\Sigma}\cC\xrightarrow{\Sigma}\cdots]\to\mathrm{Sptr}(\cC)\in\mathrm{Pr}^{L}
\]
以及极限图
\[
[\cC\xleftarrow{\Omega}\cC\xleftarrow{\Omega}\cdots]\leftarrow\mathrm{Sptr}(\cC)\in\mathrm{Pr}^{R}.
\]
于是得到伴随
\[
\Sigma^{\infty-k}:\cC\rightleftarrows\mathrm{Sptr}(\cC):\Omega^{\infty-k},
\]
其中 $\Sigma^{\infty-k}$ 由对第 $(k+1)$ 项的求值态射给出。

\begin{example}\label{L14-4.2}
对象
\[
\mathbb{S}\defeq\Sigma^{\infty}S^{0}\in\mathrm{Sptr}
\]
称为\textbf{球谱}。它在同伦代数中扮演 $\Z$ 的角色。
\end{example}

\begin{example}\label{L14-4.3}
设 $A$ 是抽象阿贝尔群。对每个 $n$，选取一个 Eilenberg--MacLane 空间 $K(A,n)$，它在同伦意义下由 $\pi_{n}K(A,n)\simeq A$ 以及 $m\neq n$ 时 $\pi_{m}K(A,n)\simeq 0$ 刻画。我们还可以选取弱同伦等价
\[
K(A,n)\xrightarrow{\ \sim\ }\Omega K(A,n+1).
\]
这些选择给出一个对象 $HA\in\mathrm{Sptr}$，它在同伦意义下良定义。我们称之为 $A$ 的\textbf{Eilenberg--MacLane 谱}。
\end{example}

\begin{remark}\label{L14-4.4}
在后续各讲中，我们将在选择空间可缩的意义下刻画 $HA$。
\end{remark}

\begin{exercise}\label{L14-4.5}
设 $E\in\mathrm{Sptr}(\cC)$。证明
\[
\colim_{k}\Sigma^{\infty-k}\Omega^{\infty-k}E\xrightarrow{\ \simeq\ }E.
\]
\end{exercise}

\begin{exercise}\label{L14-4.6}
设 $\cC$ 是紧生成的。证明对任意 $X\in\cC$ 与 $j\ge 0$，
\[
\colim_{k\ge j}\Omega^{k-j}\Sigma^{k}X\xrightarrow{\ \simeq\ }\Omega^{\infty-j}\Sigma^{\infty}X.
\]
由此推出：若 $X\in\cC$ 是紧对象，则对任意 $Y\in\cC$，有
\[
\operatorname{Maps}_{\mathrm{Sptr}(\cC)}(\Sigma^{\infty}X,\Sigma^{\infty}Y)\simeq\colim_{k}\operatorname{Maps}_{\cC}(\Sigma^{k}X,\Sigma^{k}Y).
\]
\end{exercise}

\begin{definition}\label{L14-4.7}
设 $E\in\mathrm{Sptr}$ 是谱。对任意 $n\in\Z$，定义 $E$ 的\textbf{第 $n$ 个同伦群}为
\[
\pi_{n}(E)\defeq\pi_{0}\operatorname{Maps}(\mathbb{S},\Omega^{n}E),
\]
其中当 $n<0$ 时 $\Omega^{n}\defeq\Sigma^{-n}$。
\end{definition}

\begin{remark}\label{L14-4.8}
$\pi_{n}(E)$ 是阿贝尔群，因为 $\mathrm{h}\mathrm{Sptr}$ 是加法范畴。
\end{remark}

\begin{exercise}\label{L14-4.9}
对 $Y\in\Spc_{\ast}$，证明
\[
\pi_{n}(\Sigma^{\infty}Y)\simeq\pi^{s}_{n}(Y).
\]
特别地，当 $n<0$ 时它消没。
\end{exercise}

\begin{remark}\label{L14-4.10}
上述练习说明，球的全部稳定同伦群都被编码为空间 $\operatorname{Maps}_{\mathrm{Sptr}}(\mathbb{S},\mathbb{S})$ 的通常同伦群。注意这个空间带有同伦连贯的乘法结构\footnote{我们还没有定义这意味着什么！}。
\end{remark}

\begin{exercise}\label{L14-4.11}
设 $E\in\mathrm{Sptr}$ 是谱。证明 $\Omega^{\infty}E\simeq\{\ast\}$ 当且仅当 $n\ge 0$ 时 $\pi_{n}E\simeq 0$。
\end{exercise}

\section{有限谱}

\begin{exercise}\label{L14-5.1}
设 $\cC\defeq\Ind(\cC_{0})$，其中 $\cC_{0}$ 是本质小的带点 $\infty$-范畴，配有有限极限与有限余极限，而 $\cC$ 是它的 $\Ind$-完备化。证明
\[
\mathrm{Sptr}(\cC)\simeq\Ind(\colim[\cC_{0}\xrightarrow{\Sigma}\cC_{0}\xrightarrow{\Sigma}\cC_{0}\cdots]),
\]
其中余极限在 $\Cat_{\infty}$ 中取。由此推出
\begin{itemize}[leftmargin=2.2em,itemsep=0.1em]
\item $\mathrm{Sptr}(\cC)$ 是紧生成的；
\item 函子 $\Omega^{\infty-k}:\mathrm{Sptr}(\cC)\to\cC$ 保持并检测小滤过余极限。
\end{itemize}
\end{exercise}

\begin{example}\label{L14-5.2}
对 $\cC=\Spc_{\ast}$，我们可以取 $\cC_{0}\defeq\Spc^{\mathrm{fin}}_{\ast}$，其中 $\Spc^{\mathrm{fin}}\subset\Spc$ 是包含 $\ast$ 且对有限余极限封闭的最小满子 $\infty$-范畴\footnote{一个对象属于 $\Spc^{\mathrm{fin}}$ 当且仅当它可以由有限 CW 复形表示。}。记
\[
\mathrm{Sptr}^{\mathrm{fin}}\defeq\colim[\Spc^{\mathrm{fin}}_{\ast}\xrightarrow{\Sigma}\Spc^{\mathrm{fin}}_{\ast}\xrightarrow{\Sigma}\Spc^{\mathrm{fin}}_{\ast}\cdots],
\]
并称之为\textbf{有限谱} $\infty$-范畴。我们得到等价
\[
\Ind(\mathrm{Sptr}^{\mathrm{fin}})\simeq\mathrm{Sptr},
\]
它使我们能把 $\mathrm{Sptr}^{\mathrm{fin}}$ 等同于 $\mathrm{Sptr}$ 的一个满子 $\infty$-范畴。
\end{example}

\begin{exercise}\label{L14-5.3}
函子
\[
\mathrm{Sptr}\xrightarrow{\Omega^{\infty}}\Spc_{\ast}\to\Spc
\]
保持小滤过余极限。
\end{exercise}

\begin{theorem}\label{L14-5.4}
我们有\footnote{这个结果众所周知。例如，见 MathOverflow 上的这个问题。} $\mathrm{Sptr}^{\mathrm{fin}}\simeq\mathrm{Sptr}^{\mathrm{cpt}}$。
\end{theorem}

\appendixsection{谱与上同调理论}

\begin{construction}\label{L14-A.1}
设 $E\in\mathrm{Sptr}$ 是谱。对任意 CW 对 $(X,Y)$，定义
\[
E^{n}(X,Y)\defeq\pi_{-n}(\operatorname{Maps}(\Sigma^{\infty}(X/Y),E)).
\]
记 $E^{n}(X)\defeq E^{n}(X,\varnothing)$。
\end{construction}

\begin{exercise}\label{L14-A.2}
对任意 CW 对 $(X,Y)$，构造长正合列
\[
\cdots\to E^{n}(X,Y)\to E^{n}(X)\to E^{n}(Y)\to E^{n+1}(X,Y)\to E^{n+1}(X)\to E^{n+1}(Y)\to\cdots.
\]
\end{exercise}

\begin{exercise}\label{L14-A.3}
给每个谱 $E$ 指派一个（广义）上同调理论（定义在 CW 对上）。当 $E\defeq HA$ 或 $\mathbb{S}$ 时，你分别得到什么？
\end{exercise}

\begin{exercise}\label{L14-A.4}
证明任何上同调理论都可（按上述意义）由某个谱表示，且该谱在同伦意义下唯一。
\end{exercise}

\begin{warning}\label{L14-A.5}
谱之间的非零态射可能诱导上同调理论之间的零自然变换。这样的映射称为\textbf{幻影映射}（phantom maps）。见 MathOverflow 上的这个问题。
\end{warning}

\begin{remark}\label{L14-A.6}
同调理论也有类似的故事。不过，这种构造用到谱上的 smash 积，而我们尚未定义它。
\end{remark}

\nitem\label{L14-A.7} \textbf{建议阅读}\quad HA.1.4.1。
